\subsection{System contributions: two measurement rounds and their cost findings}
\label{sec:system-cost}

Two rounds of work on this codebase, both run before this paper's own RQ1 grid and
reported here as system (not RQ1) contributions, measured where the hub's token cost
goes and how seats review each other's work. Both are dated 2026-09-18 and are recorded
in full in \texttt{docs/token-round-0918.md} and \texttt{docs/review-round-0918.md}; every
number in this subsection is copied from those two committed documents or from
\texttt{paper/budget.md}'s independent re-verification of the same rooms' launcher
telemetry, not re-derived from the underlying \texttt{swarms/*/result.json} files
directly, since those run directories are not committed to this repository (they are
local, machine-side artifacts; see the reproducibility note at the end of this
subsection). Every dollar and token figure below is therefore also regime-scoped to
2026-09-18: \texttt{paper/amendments.md} (2026-09-19 entry) records that the model
served under the \texttt{sonnet} alias (recorded model id \texttt{claude-sonnet-5}, not
a dated snapshot) changed its thinking-token budget mid-study on 2026-09-19, with no
code, flag, or reported-model-id change, moving a
single \texttt{stamp-interpreter} attempt from about 1.6K output tokens/\$0.058 to about
10K output tokens/\$0.18--0.26, and shifting arm-C per-seat output tokens from a 4K--9K
range to 18K--25K. That shift postdates both rounds reported here and does not alter the
numbers below, but it is the reason this paper states a specific date next to every cost
figure and warns against reading any \$/token ratio in this subsection as a
provider-stable constant: the same room re-run today would very likely cost more, for a
reason that has nothing to do with the hub mechanics under study.

\paragraph{The token round.} On 2026-09-17 the maintainer's OpenRouter account was
exhausted after what \texttt{docs/token-round-0918.md} records as ``roughly 2 billion
tokens'' that day; no finer-grained committed breakdown of that day's usage exists, so
this paper reports the figure at the granularity the source document gives it and does
not invent a per-room decomposition for it. Seats moved to Claude Code on Sonnet, capped
at 15 agents, and a research room measured where the tokens were going before any fix
was built. One instrumented seat ran 202 steps and processed 14,076,892 prompt tokens
against 35,298 completion tokens (a $\sim$400:1 ratio); the average of 69,687 prompt
tokens per step matches the hub's 240,000-character trim ceiling almost exactly, meaning
the seat was, on average, resending close to its full allowed context on every turn. A
fixed system-prompt-and-tool-schema tax of 5,000--6,900 tokens per turn was measured
independently three times. \texttt{wait\_for\_messages} calls were 43.8\% of that seat's
total tool calls, and idle-or-no-op calls were 39\% of all calls across the room. In one
40-minute run, one shared document was read 86--89 times by 24--25 seats. Seats given no
cap ran 203--411 steps and accumulated 10.9--18.9M prompt tokens each; seats that handed
off proactively stopped at 90--130 steps and 3--6M tokens. Across Claude Code seats
specifically, session transcripts showed a 97.0\% cache-read rate and, despite that,
still a 174:1 ratio of billed input to output tokens: caching in this system cuts price,
not the number of tokens a seat processes turn over turn.

Four rooms built and shipped a seven-item plan from that research (\texttt{docs/token-round-0918.md}):
a free re-poll loop that breaks only on an actionable result; the opt-in
\texttt{wait\_for\_messages(hold\_until\_actionable=true)} mode that holds a seat's poll
in the hub until something it must act on arrives; proactive handoff triggered on step
count, prompt-token count, and context fraction; first-reader read digests plus one
sentence of quiet-messaging guidance in the flat-run prompts; an OpenRouter sticky
session id with recorded cache fields (mock-tested only); a session-scoped tool surface
(\texttt{tool\_scope}); and an opt-in, default-off cache-stable checkpoint-trim flag,
withheld by default because summarization carries a documented answer-quality risk
\citep{Cim2026} and the oracle harness had no provider credit to test it against at the
time. Model routing across cost tiers was considered and not built, for lack of any
local measurement to size it (as distinct from the per-query cascades in
\citet{Chen2023} and \citet{Moslem2026}, which this system does not implement).

The same day produced one same-shape, same-day comparison, reported as a direction, not
a controlled result (one run per arm, different tasks): a ten-seat, read-only room on the
old build (the morning research room) against a ten-seat, read-only room on the new
build with seats told to call \texttt{hold\_until\_actionable=true}. Table~\ref{tab:token-round-first-measurement}
reproduces \texttt{docs/token-round-0918.md}'s own comparison, computed from Claude
session transcripts (\texttt{scripts/claude-room-usage.py}, deduplicated by message id) --
a measurement that depends on local, uncommitted transcript files and cannot be
regenerated from this repository's committed artifacts alone (see the note at the end of
this subsection).

\begin{table}[h]
\centering
\small
\begin{tabular}{lrr}
\toprule
 & Morning research room (old build) & Review room (new build, held waits) \\
\midrule
Worker seats & 10 & 9 \\
Model turns & 805 & 451 \\
Tokens processed (input + cache) & 115.6M & 53.6M \\
Per seat & 11.6M & 6.0M \\
Per turn & 144K & 119K \\
\texttt{wait\_for\_messages} calls & $\sim$513, none held & 42, all held \\
Messages in the room & 191 & 110 \\
Minutes to conclusion & 10.5 & 12.6 \\
\bottomrule
\end{tabular}
\caption{First same-day measurement of \texttt{hold\_until\_actionable}, from Claude
session transcripts (\texttt{docs/token-round-0918.md}, ``First measurement''). Tokens
processed fell by about half, almost entirely because held waits removed model turns;
the two rooms ran different tasks and this is one run per arm, so it is reported as a
direction only, not an effect estimate.}
\label{tab:token-round-first-measurement}
\end{table}

The launcher's own \texttt{result.json} telemetry for the review room agrees with the
transcript-derived figures: \texttt{usage.coverage} is \texttt{"complete"} for all 10
seats, with 60.1M cache-read tokens, 1.6M cache-creation tokens, 377K output tokens, and
a total API-list-price cost of \$22.34 (\texttt{docs/token-round-0918.md}; re-confirmed
directly against that room's \texttt{result.json} in \texttt{paper/budget.md}).

\paragraph{The review round.} A second round of work addressed a different gap: outside
the hub, the maintainer's own pipeline is prompt, prior-art agent, plan, build, and a
devkit review; inside a chatroom room, seats had been assigning each other work and
committing without any of that review step. A read-only research room found that the
hub's \texttt{verifiedBy()} check had been validating only that a \texttt{verify/*} board
entry existed, never what it said: 3 of 37 verify entries in one sampled run explicitly
called their own work PARTIAL, BLOCKED, or NOT GREEN and still satisfied the gate.
Volunteered review load was also unevenly concentrated: 5 of 8 verify entries in that
sample were written by one seat. Separately, commit \texttt{15d3717} had frozen new
quality baselines in the same commit as the codebase growth those baselines exist to
excuse. The same research found that a set of real, design-level defects from an earlier
morning's work was caught only by a seat reading the diff by hand while the automated
regression suite stayed green throughout: 0 of 3 real defects found that day were
devkit-shaped, and 0 of 58 sampled verify entries mentioned a size or fan-out finding --
peer review and the devkit's automated checks were catching near-disjoint classes of
failure, so neither substitutes for the other. The devkit's own deterministic bundle
costs about 2.4 seconds and no tokens per commit, but its duplication check had not
actually run on any of the 12 commits sampled that day, and its real cost is the
author's own retry turns rather than devkit's own runtime.

Three rooms built on this finding (\texttt{docs/review-round-0918.md}): a
\texttt{verify/*} entry now counts only with a machine-parseable JSON first line
naming the exact proposal id with \texttt{exit\_code:0}, exactly the gate this paper's
own methodology (\texttt{paper/protocol.md} \S9, this room's own rules) requires of
every verification claimed within it; the hub assigns a reviewer automatically when a
\texttt{claim/*} area is first taken (least-recently-verifying active non-owner,
never the same connection), tells both parties, and prefers that reviewer's entry over
anyone else's while they remain active; and a commit touching the devkit's own frozen
baselines now requires a same-gate \texttt{verify/baseline-*} entry from a different
seat, failing closed if the hub is unreachable. The same round also landed the ``lean
Claude seats'' launch path (\texttt{src/claude-args.ts}): every \texttt{claude -p} seat
is now launched with its \texttt{--tools} list restricted to that role's built-ins,
\texttt{--disable-slash-commands}, \texttt{--setting-sources project}, and
\texttt{--exclude-dynamic-system-prompt-sections}, opt-outable per seat via
\texttt{--claude-full}. A first-turn context probe on Haiku fell from 27.2K to 20.7K
tokens ($\sim$24\% reduction, one probe per combination, with 2--3K of measurement
noise on top).

Table~\ref{tab:review-round-cost} reproduces \texttt{docs/review-round-0918.md}'s own
``Cost of the day, worker seats'' table in full: seven rooms run that day, from
session-transcript-derived turn and token counts, with API-list-price dollar cost where
the launcher's own \texttt{result.json} telemetry had \texttt{coverage: "complete"}.

\begin{table}[h]
\centering
\small
\begin{tabular}{lrrrr}
\toprule
Room & Seats & Turns & Tokens processed & Cost (API price) \\
\midrule
Token research (morning) & 10 & 805 & 115.6M & n/a \\
Telemetry & 2 & 237 & 38.2M & n/a \\
Quorum and leave & 3 & 517 & 103.1M & n/a \\
Token build & 7 & 1{,}324 & 272.8M & n/a \\
Review research (held waits) & 9 & 451 & 53.6M & \$22.34 \\
Review build & 5 & 960 & 233.3M & \$70.97 \\
Lean seat build & 4 & 385 & 64.7M & \$21.99 \\
\bottomrule
\end{tabular}
\caption{\texttt{docs/review-round-0918.md}'s cost table, session-transcript-derived
(worker seats only; the review-build room's \texttt{result.json} reports 6 seats
because it also counts the verifier seat, which the transcript-derived worker count
excludes -- \texttt{paper/budget.md} flags this population mismatch explicitly). About
880M tokens across the day, 98--99\% of it cache reads. Held waits roughly halved a
research room's token count (Table~\ref{tab:token-round-first-measurement}); build rooms
did not get cheaper the same way, because a long Claude Code session's context only
grows and there is, as of this writing, no handoff or compaction mechanism for Claude
seats specifically.}
\label{tab:review-round-cost}
\end{table}

\paragraph{Cache-read accounting, explained once here.} Every cost figure in this
subsection, and every cost figure elsewhere in this paper that touches a Claude seat,
is dominated by cache-read tokens billed at a fraction of the fresh-input rate: on the
Claude Sonnet~5 rate card fetched for \texttt{paper/budget.md}
(\url{https://platform.claude.com/docs/en/about-claude/pricing}), base input is
\$2.00 per million tokens (MTok) and a cache-read hit is \$0.20/MTok -- exactly one
tenth. This is why a room that resends nearly its entire transcript on every turn (the
token round's own finding: 69,687 prompt tokens per step, matching the 240,000-character
trim ceiling) costs roughly a tenth of what its raw token count would suggest at the
base input rate, not the full base-rate price: the 233M-token review-build session
(Table~\ref{tab:review-round-cost}) had about 280M cache-read tokens, which alone would
cost about \$560 at the \$2.00/MTok base input rate, yet its measured cost was \$70.97,
because those tokens were cache-read hits billed at \$0.20/MTok, not fresh input.
Concretely, for that same session \texttt{paper/budget.md} records
\texttt{usage.input\_tokens: 2{,}354}, \texttt{usage.cache\_read\_input\_tokens:
279{,}597{,}966}, \texttt{usage.cache\_creation\_input\_tokens: 2{,}108{,}427}, and
\texttt{usage.output\_tokens: 657{,}972} against a measured \texttt{cost\_usd} of
\$70.9735702 -- a raw repricing of these four fields at the published per-MTok rates
undershoots the CLI-reported cost by 4.5\% for this session (and by 4.5--11.3\% across
the three rooms \texttt{paper/budget.md} checked), a gap that document attributes,
tentatively and without full attribution, to 1-hour cache writes (a \$4/MTok premium
tier on Sonnet~5, above the plain 5-minute cache-write rate of \$2.50/MTok) or a
data-residency multiplier not broken out in \texttt{result.json}'s fields. This paper
therefore always reports the CLI's own measured \texttt{cost\_usd} as authoritative
where it is available, and treats a from-scratch repricing of the four token fields as a
sanity bound accurate to 4.5--11.3\%, never as an exact reconstruction.

\paragraph{A note on what is, and is not, reproducible from this repository.} The two
per-room dollar figures and token-field breakdowns quoted above (Table~\ref{tab:review-round-cost}'s
\$22.34/\$70.97/\$21.99 column, and this paragraph's four \texttt{usage} fields) come
from each room's own \texttt{swarms/<room>/result.json}, which is the launcher's
committed-schema run artifact (\texttt{src/result.ts}) but whose individual instances
are not committed to this repository (the \texttt{swarms/} directory is gitignored); a
reader can regenerate the same class of number for a fresh run by inspecting that run's
own \texttt{result.json}, but cannot recompute \emph{these specific historical rooms'}
numbers without those local files, which is why this paper cites them only through the
two committed documents (\texttt{docs/token-round-0918.md}, \texttt{docs/review-round-0918.md})
and \texttt{paper/budget.md}'s own re-verification, rather than asserting a direct
read of an artifact this repository does not ship. The turn and token counts in both
tables above (``Tokens processed'', ``Model turns'') are, separately, derived from local
Claude Code session transcripts via \texttt{scripts/claude-room-usage.py}, an accounting
this project's own system description (\texttt{paper/system.md} \S3) is explicit is
\emph{not} the same quantity as \texttt{result.json}'s \texttt{usage} rollup (which is a
run-level aggregate, not a per-seat breakdown, for Claude seats) and cannot be
regenerated from this repository's committed artifacts alone.
